\documentstyle[12pt,fullpage]{article}

\textheight 24cm

\begin{document}

\title{Computational aspects of the N-body problem }
\author{ 
ILIAS KOTSIREAS 
\\LIP6, CalFor, Case 168, Universit\'e Paris 6, 
\\4, place Jussieu, 75252 Paris Cedex 05, France
\\E-mail\,: Ilias.Kotsireas@lip6.fr
\\web\,: http://posso.lip6.fr/$ \ \tilde{\left. \right.} \ $ilias/}
\date{January 1999}
\maketitle

The study of central configurations 
in the N-body problem of Celestial Mechanics
offers great computational challenges
that can be successfully settled using 
Computer Algebra methods and techniques.  \par
The general aim is to search and classify
all possible types of central configurations.
This is especially interesting 
for the planar newtonian 4-body problem with equal masses and
for the spatial newtonian 5-body problem with equal masses
as well as for various types of corresponding problems 
with non equal masses. \par  
The fundamental fact that allows us to apply Computer Algebra methods to
the problem of central configurations is the description of
the equations of central configurations by a system of polynomial 
equations. A theoretical tool that is used to simplify 
the problem is the demonstration of the symmetry theorems.
The original polynomial equations are processed using the
matrix formalism in order to obtain an equivalent but easier
to solve system. The algorithmic tools used are resultants, Gr\"obner bases,
Triangular sets and some Invariant theory of Finite Groups.
The Computer Algebra programs used are Axiom, Macsyma, Maple, Pari,
NTL, Gb, FGb, Magma, MPS, RealSolving and RS. \par
On the bare computational level the study of central configurations
gives rise to zero-dimensional polynomial systems in the case of equal masses.
Some of these systems are easy to solve but others are inaccessible.
In the case of non equal masses one obtains parametric polynomial systems. \par
Another interesting computational aspect of the problem is the use of
quadratic forms to interpret the determinental identities that arise
in the proofs of the symmetry theorems. The study of the corresponding 
isotropic cones might lead to easier proofs. Since the number of
variables of these quadratic forms is more than $10$, 
a Maple package called {\tt QF} has been developed in order 
to study them. \par
The study of central configurations gives rise to a wide variety of
challenging problems that are suitable for Computer Algebra.
The key notion is a better understanding of the symmetries.   

\end{document}
